\documentclass[journal]{IEEEtran}
\usepackage{amsmath,amssymb,mathtools}
\usepackage{amsthm}
\usepackage{booktabs,array}
\usepackage{enumitem}
\usepackage{microtype}
\usepackage[hidelinks]{hyperref}
\newtheoremstyle{plainn}{}{}{\itshape}{}{\bfseries}{.}{.5em}%
  {\thmname{#1}\thmnumber{ #2}\thmnote{ #3}}
\theoremstyle{plainn}
\newtheorem*{lem}{Lemma}
\newtheorem*{thm}{Theorem}
\newtheorem*{prop}{Proposition}
\newtheorem*{cor}{Corollary}
\newcommand{\rd}{\mathrm{d}}

\begin{document}
\title{Parameterized algorithms for functional observer feasibility}
\author{}
\maketitle

\begin{abstract}
We study the computational complexity of deciding whether a prescribed-order functional observer exists for a linear system with rational input data and a fixed observer row block. Building on the existing algebraic characterization of functional-observer feasibility, we develop two parameterized decision algorithms. First, we show that feasibility is decidable in deterministic polynomial time for every fixed observer order $q$, with a bit-complexity bound of $\mathcal L^{O(q^2)}$, where $\mathcal L$ is the input bit length. Second, we establish fixed-parameter tractability with respect to the effective free-row dimension
\[
\kappa_{\rm obs}=r\operatorname{rank}[N\ AN],
\]
obtaining a bound $f(\kappa_{\rm obs})\operatorname{poly}(\mathcal L)$; the algorithm combines a direct parameterization of the attainable free-row pairs with a finite algebraic certificate for the spectral rank condition. We also close the boundary case of no free observer rows by an exact polynomial-time test based on polynomial left inverses, and we establish membership of the prescribed-order decision problem, and of its finite threshold version, in the existential theory of the reals. These results are parameterized upper bounds; they do not resolve the unrestricted complexity of functional-observer feasibility.
\end{abstract}

\begin{IEEEkeywords}
functional observer; minimum order; parameterized complexity; real algebraic geometry; polynomial matrix; fixed-parameter tractability.
\end{IEEEkeywords}

\section{Introduction}\label{sec:intro}
\subsection{The problem}\label{sec:problem}
Let $A\in\mathbb R^{n\times n}$, $C\in\mathbb R^{p\times n}$, and $L_0\in\mathbb R^{r_0\times n}$ be given. For an integer $q\ge r_0$, let

\[
L=\begin{bmatrix}L_0\\ R\end{bmatrix}\in\mathbb R^{q\times n}
\]
and impose

\[
(5a)\quad \operatorname{rank}\begin{bmatrix}LA\\ L\\ CA\\ C\end{bmatrix}
=\operatorname{rank}\begin{bmatrix}L\\ CA\\ C\end{bmatrix},
\]
\[
(5b)\quad \operatorname{rank}\begin{bmatrix}\lambda L-LA\\ CA\\ C\end{bmatrix}
=\operatorname{rank}\begin{bmatrix}L\\ CA\\ C\end{bmatrix}
\quad\forall\lambda\in\mathbb C.
\]
These are the Fernando--Trinh--Jennings conditions for a $q$-order functional observer \cite{fer10}. The decision version $\mathrm{DEC}(q)$ asks whether such an $R$ exists for a given $q$; the optimization version asks for the minimum feasible $q$, denoted $q_{\min}$. The name $q1$ denotes the complexity of the \emph{decision} problem $\mathrm{DEC}$; see Section~\ref{sec:io} for the exact convention.

The characterization problem is settled: Fernando \cite{fer26a,fer26b} gives the order spectrum, per-order necessary and sufficient conditions, the minimum order, and a construction. The complexity of the resulting existential problem remains open. The present paper does not settle it; it identifies two structural parameters under which the problem is tractable.

\subsection{Contributions}\label{sec:contrib}
\begin{enumerate}
\item \textbf{Compression and pencil criterion.} We reduce $(5a)(5b)$ to a pair of matrices $H=LN$, $H_A=LAN$, with $N$ a basis of $\ker[C;CA]$; $(5a)$ becomes a row-space inclusion and $(5b)$ becomes the non-vanishing of the maximal minors of the pencil $\lambda H-H_A$.
\item \textbf{Fixed observer order.} For every fixed $q$, we give a polynomial-time decision algorithm, with a bound $\mathcal L^{O(q^2)}$. The algorithm uses a rank factorization $H=XB$, the substitution $G=FX$, a parametric linear system in the remaining entries of $B$, rank-consistency elimination, and a $2\omega$-variable spectral certificate.
\item \textbf{Fixed effective dimension.} We give an FPT algorithm parameterized by $\kappa_{\rm obs}=r\operatorname{rank}[N\ AN]$, with bound $f(\kappa_{\rm obs})\operatorname{poly}(\mathcal L)$. The bound covers every $r\ge0$; the boundary value $r=0$, where the parameter degenerates, is closed separately (item 4).
\item \textbf{The zero-free-row case.} For $r=0$ the two conditions reduce to a constant-rank test and the constancy of the rank of a degree-one pencil $\lambda X-G$. We show that the latter is equivalent to the consistency of a finite rational linear system with $O(q\omega^2)$ unknowns and $2\omega^3$ equations, derived from the polynomial-left-inverse degree bound $\deg S\le2\omega-2$. This gives a deterministic exact $\operatorname{poly}(q,\omega)$ algorithm, with an explicit witness, and removes the only parameter value not covered by item 3.
\item \textbf{Real-algebraic upper bound and combination.} We record $\mathrm{DEC}(q)\in\exists\mathbb R$ (hence $\mathrm{MINDEC}(Q)\in\exists\mathbb R$) and the combined bound of Theorems 1 and 2 (Corollary 4).
\end{enumerate}
The two algorithms control different sources of complexity:

\begin{table*}[!t]
\centering
\caption{The two parameterized algorithms and what each controls.}\label{tab:1}
\small
\begin{tabular}{p{0.300\textwidth}p{0.300\textwidth}p{0.300\textwidth}}
\toprule
Feature & Fixed $q$ & Fixed $\kappa_{\rm obs}$ \\
\midrule
Parameter & $q$ & $\kappa_{\rm obs}=r\,d_0$ \\
Main reduction & pivot factorization $H=XB$ & direct free-row parametrization \\
Advantage & polynomial time for fixed order & polynomial exponent independent of the parameter \\
Limitation & exponent depends on $q$ & relies on real-feasibility bounds \\
Key structural fact & low-dimensional pencil variables & low-dimensional free-row space \\
\bottomrule
\end{tabular}
\end{table*}
The first result provides a complementary XP-type upper bound.

\subsection{What is not claimed}\label{sec:notclaim}
We do not claim that $q1\in\mathrm P$, nor that it is $\exists\mathbb R$-complete, nor any hardness result. The two upper bounds are compatible with all three outcomes: a polynomial algorithm, $\exists\mathbb R$-completeness, or an intermediate parameterized classification. Hardness and gadget constructions are outside the scope of this paper.

\subsection{Relation to the characterization literature}\label{sec:rel}
Fernando \cite{fer26a,fer26b} characterizes the minimum-order functional observer: the order spectrum $\Sigma_{FO}\subseteq[q_0,n_0-p]$, per-order necessary and sufficient conditions, the minimum $\nu_{\min}=\min\Sigma_{FO}$, and a construction; a more recent observer-design framework is \cite{dar25}. The present paper takes that characterization as given and asks a different question, namely the \emph{computational} cost of deciding feasibility at a prescribed order. Two consequences matter for reading what follows.

\begin{enumerate}
\item None of the structural statements of \cite{fer26a,fer26b} is re-derived here. Conditions $(5a)$, $(5b)$ are used as the definition of feasibility, and the compression of Lemma A is the standard reduction to the restricted pencil; where an ingredient is classical (rank factorization, Gram matrices, univariate B\'ezout), this is stated at the point of use.
\item The results are upper bounds. They do not contradict, and are compatible with, any of the three possible outcomes for $q1$.
\end{enumerate}
The division of labour is summarized below. The right-hand column is a statement about \cite{fer26a,fer26b} only; it is not an exhaustive claim about the literature, and Section~\ref{sec:paramsize} records where a comparison is still owed.

\begin{table*}[!t]
\centering
\caption{Division of labour with the characterization literature.}\label{tab:2}
\small
\begin{tabular}{p{0.300\textwidth}p{0.300\textwidth}p{0.300\textwidth}}
\toprule
Question & \cite{fer26a,fer26b} & This paper \\
\midrule
Which orders are feasible? & yes (order spectrum) & uses it \\
Algebraic conditions at a prescribed order & yes & uses them \\
Spectral pencil condition & yes (pencil form, Thm 4/5) & uses it \\
Polynomial time for fixed order? & not established there & Theorem 1 \\
FPT in a problem-specific parameter? & not established there & Theorem 2 \\
Zero-free-row case by a finite linear system? & not established there & Proposition 3c \\
\bottomrule
\end{tabular}
\end{table*}
The contribution is therefore not a new characterization. It is a reduction that makes the decision problem finite-dimensional, together with two parameterized consequences of that reduction. Accordingly, the paper addresses the computational complexity of prescribed-order feasibility rather than providing another characterization of the feasible observer-order spectrum.

\section{Setting and notation}\label{sec:setting}
\subsection{Input, output, and complexity conventions}\label{sec:io}
The input $A,C,L_0$ has rational entries encoded in binary; $\mathcal L$ denotes the total bit length. We use the real observer-design convention:

\[
A,C,L_0\in\mathbb Q,\qquad L=[L_0;R]\in\mathbb R^{q\times n},\qquad
\lambda\in\mathbb C.
\]
Ranks are computed over $\mathbb C$ (for real matrices this agrees with the rank over $\mathbb R$). The spectral variable $\lambda$ remains complex. If one allows a complex $R$, realification doubles the free coordinates and replaces every transpose by the conjugate transpose; the FPT statement is unchanged because $f(2\kappa_{\rm obs})=\widetilde f(\kappa_{\rm obs})$. $q\le n$ throughout.

Formally, the decision language is

\[
\begin{aligned}
\mathrm{DEC}(q)=\Bigl\{(A,C,L_0,q):\ &\exists R\in\mathbb R^{(q-r_0)\times n}\\[-2pt]
&\text{such that }(5a),(5b)\text{ hold}\Bigr\}.
\end{aligned}
\]
The optimization problem $\mathrm{OPT}$ asks for the least $q$ with $(A,C,L_0,q)\in\mathrm{DEC}(q)$; $\mathrm{FIND}$ asks for a witness $R$ at the optimum. This paper studies the \emph{decision} problem: throughout, $q1$ denotes the complexity of deciding $\mathrm{DEC}$, and $\mathrm{OPT}$, $\mathrm{FIND}$ are function problems that are not addressed here. We do not use monotonicity of $\mathrm{DEC}(q)$ in $q$, and we do not claim that $\mathrm{OPT}$ and $\mathrm{FIND}$ have the same complexity.

All complexity statements below refer to the decision problem. Membership in $\exists\mathbb R$ is proved for $\mathrm{DEC}(q)$ at each fixed $q$, and therefore for the threshold problem

\[
\mathrm{MINDEC}(Q)=\bigvee_{q=r_0}^{\min(Q,n)}\mathrm{DEC}(q)
\]
(a finite disjunction, so the encoding size is additive). Claims of the form "$q1\in\exists\mathbb R$" are to be read in this sense; $\mathrm{OPT}$ itself is a function problem and is not a language in $\exists\mathbb R$.

\subsection{Observer compression}\label{sec:compress}
Let

\[
K_0=\ker\begin{bmatrix}C\\ CA\end{bmatrix},\qquad k=\dim K_0,
\]
and let $N\in\mathbb R^{n\times k}$ be a basis. Define

\[
H=LN\in\mathbb R^{q\times k},\qquad H_A=LAN\in\mathbb R^{q\times k},\qquad
\omega=\operatorname{rank}H.
\]
\subsection{The free-row space}\label{sec:free}
For a free row $r_i\in\mathbb R^{1\times n}$,

\[
(H_i,(H_A)_i)=r_i[N\ AN].
\]
Hence all free row pairs lie in

\[
\mathcal T=\operatorname{row}[N\ AN]\subseteq\mathbb R^{2k},
\qquad d_0=\dim\mathcal T=\operatorname{rank}[N\ AN].
\]
\subsection{Rank conventions}\label{sec:rankconv}
For a matrix $M$, $\operatorname{rank}M$ is the dimension of its row space. For a pencil $P(\lambda)$, its $\omega\times\omega$ minors are the maximal minors when $\operatorname{rank}P(\lambda)=\omega$.

\section{Lemma A: compression and pencil criterion}\label{sec:lemA}
\begin{lem}[A] Let $N$ be a basis of $K_0$, $H=LN$, $H_A=LAN$, and $\omega=\operatorname{rank}H$.

(i) $(5a)$ holds if and only if

\[
\operatorname{row}H_A\subseteq\operatorname{row}H.
\]
(ii) If $(5a)$ holds, there exists $F\in\mathbb R^{q\times q}$ such that

\[
H_A=FH.
\]
(iii) $(5b)$ holds if and only if the $\omega\times\omega$ maximal minors of

\[
P(\lambda)=\lambda H-H_A
\]
have no common complex root; equivalently, their gcd in $\mathbb C[\lambda]$ is a nonzero constant.

(iv) If $(5a)$ and $(5b)$ hold, then

\[
\operatorname{rank}H_A=\operatorname{rank}H.
\]
\end{lem}
\emph{Proof.} (i) is the row-space form of the rank equality defining $(5a)$: $\operatorname{rank}[LA;L;CA;C]=\operatorname{rank}[L;CA;C]$ is equivalent to $\operatorname{row}(LA)\subseteq\operatorname{row}[L;CA;C]$. Right multiplication by $N$ kills $CA$ and $C$, so $\operatorname{row}H_A\subseteq\operatorname{row}H$. Conversely, if $H_A=\Lambda H$, then $(LA-\Lambda L)N=0$, so every row of $LA-\Lambda L$ lies in $K_0^{\perp}=\operatorname{row}[C;CA]$, which is contained in $\operatorname{row}[L;CA;C]$; hence $(5a)$. (ii) expresses each row of $H_A$ as a linear combination of the rows of $H$; collect the coefficients in $F$. (iii) is the standard maximal-minor criterion for a rectangular pencil to have constant rank $\omega$: a common root of all maximal minors is exactly a $\lambda$ at which the rank drops. (iv) follows from (iii) at $\lambda=0$: $\operatorname{rank}H_A=\operatorname{rank}H$.

Both rank statements have the same one-line justification. Put $B_0=[CA;C]$ and let $N$ be a basis of $\ker B_0$; then for every matrix $M$,

\[
\operatorname{rank}\begin{bmatrix}M\\B_0\end{bmatrix}
=\operatorname{rank}B_0+\operatorname{rank}(MN).
\]
With $M=[LA;L]$ this gives $(5a)\iff\operatorname{rank}[H_A;H]=\operatorname{rank}H$, and with $M=\lambda L-LA$ it gives

\[
(5b)_\lambda\iff\operatorname{rank}(\lambda H-H_A)=\operatorname{rank}H .
\]
The second identity is what makes the spectral condition a \emph{constant-rank pencil} condition rather than a mere row-space inclusion: a rank equality at every $\lambda$, not an inclusion.

\emph{Interpretation.} Under $(5a)$ one has $H_A=FH$, hence

\[
\lambda H-H_A=(\lambda I-F)H .
\]
So $(5b)$ reads: at every finite complex $\lambda$, the map $\lambda I-F$ destroys no direction of $\operatorname{row}H$. This is the form in which the condition is used below.

\emph{Status.} The abstract compression criterion is Lean-verified in \texttt{Q1Paper.ker\_\allowbreak{}compressed\_\allowbreak{}iff} (\texttt{Compression.lean}). The matrix-level pencil/gcd bridge remains paper-level and is part of Appendix A.

\section{Lemma B: free-row parametrization}\label{sec:lemB}
\begin{lem}[B] With $\mathcal T=\operatorname{row}[N\ AN]$ and $d_0=\operatorname{rank}[N\ AN]$, every free row pair lies in $\mathcal T$. Choosing a rational basis of $\mathcal T$, each free row is represented by exactly $d_0$ real coordinates, and all free rows by

\[
\kappa_{\rm obs}=r d_0
\]
coordinates. Moreover,

\[
k\le d_0\le\kappa_{\rm obs}
\]
when $r\ge1$.

\end{lem}
\emph{Proof.} The identity $(H_i,(H_A)_i)=r_i[N\ AN]$ gives membership. The dimension count is the rank-nullity theorem for $r_i\mapsto(r_iN,r_iAN)$, whose image is $\mathcal T$. Since $N$ has full column rank, its columns are independent in $[N\ AN]$, hence $k\le d_0$; and $d_0\le rd_0=\kappa_{\rm obs}$ for $r\ge1$.

\emph{Basis invariance.} Replacing $N$ by $NQ$ with $Q$ nonsingular changes the coordinates of $\mathcal T$ but neither $\mathcal T$ nor $d_0=\dim\mathcal T$; likewise the coordinates $t$ change by $Q$, so the number of parameters does not depend on the choice of basis.

\emph{Example.} Take $n=4$, $C=e_1^{\mathsf T}$, and $A$ the cyclic shift $e_1\mapsto e_2\mapsto e_3\mapsto e_4\mapsto e_1$. Then $K_0=\ker[C;CA]=\operatorname{span}\{e_3,e_4\}$, so $k=2$, and $AN=\operatorname{span}\{e_4,e_1\}$, so $d_0=\dim\operatorname{span}\{e_1,e_3,e_4\}=3$. With one free row, $\kappa_{\rm obs}=3$: the four-dimensional row $r_1$ contributes $3$ coordinates, not $4$. The compressed pair $(H,H_A)$ has two columns, and the certificate of Lemma D adds $2\omega$ variables, so the formula of Theorem 2 has $3+2\omega$ real variables here.

\emph{Lean status.} The free-pair space and dimension count are verified in \texttt{Q1Paper.FreeRows} (\texttt{freePairSpace}, \texttt{finrank\_\allowbreak{}freePairSpace\_\allowbreak{}add\_\allowbreak{}ker}, \texttt{kappaObs\_\allowbreak{}eq\_\allowbreak{}mul\_\allowbreak{}freePairSpace}); the matrix-level $\mathcal T$ identification with $\operatorname{row}[N\ AN]$ is paper-level.

\section{Lemma C: rank factorization and consistency}\label{sec:lemC}
\subsection{Pivot factorization}\label{sec:pivot}
Let $\omega>0$. For a rank-$\omega$ matrix $H$, choose a pivot set $J\subseteq[k]$, $|J|=\omega$, such that

\[
X=H_{:,J}\in\mathbb R^{q\times\omega}
\]
has full column rank. Then

\[
H=XB,\qquad B\in\mathbb C^{\omega\times k},\qquad B_{:,J}=I_\omega.
\]
\subsection{The $G=FX$ substitution}\label{sec:gsub}
\emph{Lean status.} The algebraic substitution is verified in \texttt{Q1Paper.G\_\allowbreak{}form\_\allowbreak{}iff}: if $X$ has a left inverse, then $H_A=F(XB)$ for some $F$ iff $H_A=GB$ for some unconstrained $G$.

If $H_A=FH$, put

\[
G=FX\in\mathbb R^{q\times\omega}.
\]
Then

\[
H_A=GB,\qquad P(\lambda)=(\lambda X-G)B.
\]
Conversely, if $H_A=GB$ and $X_LX=I_\omega$, then $F=GX_L$ satisfies $FH=H_A$.

\subsection{Parametric linear system}\label{sec:param}
Let $D$ be a rational matrix with $\ker D=\mathcal T$. The fixed rows satisfy

\[
X_iB=h_i^0,\qquad G_iB=g_i^0,
\]
and every free row satisfies

\[
D\begin{bmatrix}(X_iB)^{\mathsf T}\\ (G_iB)^{\mathsf T}\end{bmatrix}=0.
\]
After substituting $B_{:,J}=I_\omega$, these equations are linear in the remaining entries $b$ of $B$:

\[
M(X,G)b=c(X,G).
\]
\subsection{Rank consistency}\label{sec:rankcons}
\emph{Lean status.} The subspace form is verified in \texttt{Q1Paper.RankConsistency}: consistency of \texttt{T b = c} is equivalent to \texttt{span\{c\} $\le$ range T}, equivalently \texttt{range T = range T $\sqcup$ span\{c\}}. The matrix Gram/rank identity is paper-level and checked by \texttt{c311}.

\begin{lem}[C] For real matrices $M,c$,

\[
\exists b:\ Mb=c
\]
if and only if

\[
\operatorname{rank}(MM^{\mathsf T})
=\operatorname{rank}(MM^{\mathsf T}+cc^{\mathsf T}).
\]
Equivalently, with $\Gamma=MM^{\mathsf T}$, $\Gamma'=MM^{\mathsf T}+cc^{\mathsf T}$,

\[
\bigvee_{t=0}^{d_1}
\left[e_t(\Gamma)>0\wedge e_{t+1}(\Gamma)=0\wedge e_{t+1}(\Gamma')=0\right],
\]
where $d_1=\min(\text{rows},\text{cols})$ of $M$ and $e_j$ is the $j$-th elementary symmetric polynomial of the eigenvalues.

\end{lem}
\emph{Proof.} $Mb=c$ is consistent iff $c\in\operatorname{im}M$. Since $\operatorname{im}\Gamma=\operatorname{im}M$ and $\operatorname{im}\Gamma'=\operatorname{im}[M\ c]$, this is equivalent to equality of the two ranks. For a positive semidefinite matrix $\Gamma$, $e_t(\Gamma)>0$ and $e_{t+1}(\Gamma)=0$ exactly when $\operatorname{rank}\Gamma=t$; the last condition forces $\operatorname{rank}\Gamma'\le t$, and $\Gamma'\succeq\Gamma$ gives equality.

\subsection{Spectral reduction}\label{sec:spectral}
Because $B$ has full row rank,

\[
\operatorname{rank}P(\lambda)=\operatorname{rank}((\lambda X-G)B)
=\operatorname{rank}(\lambda X-G).
\]
Thus $(5b)$ is equivalent to

\[
\operatorname{rank}(\lambda X-G)=\omega\quad\forall\lambda\in\mathbb C.
\]
\section{Lemma D: spectral certificate}\label{sec:lemD}
\subsection{Bezout degree bound}\label{sec:bezout}
\begin{lem}[D1] Let $\omega\ge1$ and let $f_1,\ldots,f_s\in\mathbb R[\lambda]$ have degree at most $\omega$ and no common complex root. Then there exist real polynomials $g_i$ with $\deg g_i<\omega$ such that

\[
\sum_i g_i f_i=1.
\]
\emph{Lean status.} The two-polynomial Euclidean-gcd form is verified in \texttt{Q1Paper.bezout\_\allowbreak{}of\_\allowbreak{}gcd\_\allowbreak{}eq\_\allowbreak{}one}; the degree bound below is paper-level.

\end{lem}
\emph{Proof.} Choose a nonzero real combination $g=\sum a_i f_i$. At each complex root $\alpha$ of $g$, the condition $h(\alpha)=0$ for $h=\sum b_i f_i$ is a proper real linear condition on $b$; avoiding the finite union of these proper subspaces gives $h$ with no common root with $g$. Then $g,h$ are coprime and the two-polynomial Bezout identity gives multipliers of degree less than $\omega$; expanding gives the claim.

\subsection{Gram certificate}\label{sec:gram}
Let $f_i$ be the maximal minors of $P(\lambda)$. Form a matrix $B$ whose rows are indexed by $(i,j)$, $0\le j<\omega$, and whose columns are indexed by $t=0,\ldots,2\omega-1$, with

\[
B_{(i,j),t}=c_{i,t-j},
\]
where $c_{i,a}$ is the coefficient of $\lambda^a$ in $f_i$. Then

\[
(5b)\iff e_0\in\operatorname{row}B
\iff\exists x\in\mathbb R^{2\omega}:\ B^{\mathsf T}Bx=e_0,
\]
with $e_0=(1,0,\ldots,0)$.

\emph{Proof.} The rows of $B$ are the coefficient vectors of $\lambda^j f_i$. Their row space is exactly the set of coefficient vectors of $\sum g_i f_i$ with $\deg g_i<\omega$ (Lemma D2 below). Lemma D1 gives a representation of $1$ exactly when $(5b)$ holds. For real $B$, $\operatorname{row}B=\operatorname{im}(B^{\mathsf T}B)$, giving the equivalent linear system.

\begin{lem}[D2] Let $f_1,\dots,f_s\in\mathbb R[\lambda]$ have degree at most $\omega$, let $c_{i,a}$ be the coefficient of $\lambda^a$ (with $c_{i,a}=0$ for $a<0$ or $a>\omega$), and let $B$ be the matrix with rows $(i,j)$, $0\le j<\omega$, columns $t=0,\dots,2\omega-1$, and $B_{(i,j),t}=c_{i,t-j}$. Then, reading vectors in the basis $1,\lambda,\dots,\lambda^{2\omega-1}$,

\[
\operatorname{row}B
=\Bigl\{\,\text{coefficient vector of }\textstyle\sum_i g_if_i
\ :\ \deg g_i<\omega\,\Bigr\}.
\]
\end{lem}
\emph{Proof.} Row $(i,j)$ of $B$ is the coefficient vector of $\lambda^jf_i$. Hence a combination with weights $y_{i,j}$ is the coefficient vector of $\sum_{i,j}y_{i,j}\lambda^jf_i=\sum_ig_if_i$ with $g_i=\sum_jy_{i,j}\lambda^j$, and $\deg g_i<\omega$; conversely every $\sum_ig_if_i$ with $\deg g_i<\omega$ arises this way. Since $\deg(\lambda^jf_i)\le\omega+j\le2\omega-1$, all coefficients fit in the displayed $2\omega$ columns, and the convention $c_{i,a}=0$ outside $0\le a\le\omega$ makes the truncation automatic. $\square$

\subsection{Cauchy--Binet compression}\label{sec:cb}
Define

\[
K(\lambda,\mu)=\sum_{|J|=\omega}
\det\left(P(\lambda)_{:,J}^{\mathsf T}P(\mu)_{:,J}\right).
\]
For each fixed column set $J$, apply Cauchy--Binet to the $\omega\times k$ matrix $P(\lambda)_{:,J}^{\mathsf T}$ and the $k\times\omega$ matrix $P(\mu)_{:,J}$:

\[
\det\!\left(P(\lambda)_{:,J}^{\mathsf T}P(\mu)_{:,J}\right)
=\sum_{|I|=\omega}f_{I,J}(\lambda)f_{I,J}(\mu),
\]
where $f_{I,J}$ is the minor on rows $I$ and columns $J$. Summing over $J$ gives

\[
K(\lambda,\mu)=\sum_{I,J}f_{I,J}(\lambda)f_{I,J}(\mu),
\]
the Gram kernel of \textbf{all} maximal minors of $P$.

Two consequences are used below.

\begin{enumerate}
\item \emph{Cost.} Computing $K$ needs only $\binom{k}{\omega}\le2^{\kappa_{\rm obs}}$ determinants of order $\omega$, one per column set $J$; the $\binom q\omega\binom k\omega$ individual minors are never enumerated.
\item \emph{Coefficient identity.} Writing $S_{a,b}$ for the coefficient of $\lambda^a\mu^b$ in $K$ and $c_{i,a}$ for the coefficient of $\lambda^a$ in $f_i$, one has $S_{a,b}=\sum_ic_{i,a}c_{i,b}$, hence, with the row indexing of Lemma D2, \[
(B^{\mathsf T}B)_{t,t'}
   =\sum_{j=0}^{\omega-1}\sum_ic_{i,t-j}c_{i,t'-j}
   =\sum_{j=0}^{\omega-1}S_{t-j,t'-j}.
\]
\end{enumerate}
\emph{Algorithmic reading.} The identity is classical (an exterior-power Gram identity; see e.g. \cite{emi99} for the B\'ezoutian/elimination background). What the algorithm uses is the compression, and it is worth writing the four steps out:

\begin{enumerate}
\item for each column set $J\subseteq[k]$, $|J|=\omega$, form the $\omega\times\omega$ pencil $P(\lambda)_{:,J}^{\mathsf T}P(\mu)_{:,J}$;
\item compute its bivariate coefficients (degree at most $\omega$ in each variable);
\item sum the coefficients over $J$;
\item recover $B^{\mathsf T}B$ from the shifted identity above.
\end{enumerate}
The cost is $\binom{k}{\omega}$ determinants of order $\omega$; the $\binom q\omega\binom k\omega$ minors of $P$ are never formed, and no hidden enumeration replaces them. In particular this is what keeps the certificate matrix parameter-sized when $q$ is large and $k$ is bounded.

\begin{table*}[!t]
\centering
\caption{Classical tools and their role in the construction.}\label{tab:3}
\small
\begin{tabular}{p{0.450\textwidth}p{0.450\textwidth}}
\toprule
Classical tool & Role in this paper \\
\midrule
Cauchy--Binet & produces the Gram-kernel identity \\
compound matrices & explain the maximal-minor structure \\
B\'ezout identities & convert spectral feasibility into a finite certificate \\
polynomial-matrix left inverses & certificate existence and the degree bound \\
Renegar's bound & the bit-complexity bound \\
\bottomrule
\end{tabular}
\end{table*}
The individual ingredients are classical; the contribution lies in their algorithmic combination, namely the use of the compression to keep the certificate parameter-sized.

\section{Theorem 1: fixed observer order}\label{sec:thm1}
\begin{thm}[1] For every fixed $q$, $q1$ is decidable in deterministic polynomial time. More precisely, the algorithm described below runs in

\[
\mathcal L^{O(q^2)}
\]
bit operations.

\emph{Algorithm.} Enumerate $\omega=0,\ldots,q$ and the pivot sets $J\subseteq\{1,\dots,k\}$, $|J|=\omega$. For each chart, form the existential real formula $\Phi_{\omega,J}(X,G,x)$ with the following constraints:

\begin{enumerate}
\item $X\in\mathbb R^{q\times\omega}$, $G\in\mathbb R^{q\times\omega}$, $x\in\mathbb R^{2\omega}$;
\item $e_\omega(X^{\mathsf T}X)>0$ (the left inverse $U$ used in the proof below is not quantified);
\item the rank-consistency formula of Lemma C for the parametric system $M(X,G)b=c(X,G)$, with $B_{:,J}=I_\omega$;
\item the fixed-row equations $X_iB=h_i^0$, $G_iB=g_i^0$ for $i\le r_0$;
\item the free-row equations $D[(X_iB)^{\mathsf T};(G_iB)^{\mathsf T}]=0$;
\item the spectral certificate $B^{\mathsf T}Bx=e_0$, where $B$ is built from the maximal minors of $\lambda X-G$.
\end{enumerate}
Run a real-feasibility algorithm on the finite disjunction $\bigvee_{\omega,J}\Phi_{\omega,J}$.

\emph{Correctness.} Suppose a feasible $L$ exists. Choose a pivot set $J$ of $H$ and the factorization $H=XB$ with $B_{:,J}=I_\omega$. By Lemma A, $H_A=FH$; put $G=FX$ and $U$ a left inverse of $X$. The fixed and free row equations hold by construction of $H=LN$, $H_A=LAN$. Lemma C gives the rank-consistency condition, and Lemma D gives the spectral certificate. Hence the formula is true.

Conversely, suppose $\Phi_{\omega,J}$ is true. Rank consistency gives $B$ with $B_{:,J}=I_\omega$; set $H=XB$, $H_A=GB$. The fixed-row equations identify the first $r_0$ rows with $L_0N,L_0AN$; the free-row equations put each remaining row pair in $\mathcal T$, so each has a preimage $r_i$ with $(H_i,(H_A)_i)=r_i[N\ AN]$. Let $L=[L_0;R]$. Then $H=LN$, $H_A=LAN$. Choose any left inverse $U$ of $X$ (it exists because $e_\omega(X^{\mathsf T}X)>0$); then $G=(GU)X$, so $H_A=GB=(GU)XB=FH$ with $F=GU$, and Lemma A gives $(5a)$. The spectral certificate and Lemma D give $(5b)$. Thus the formula is equivalent to the existence of a feasible $L$.

\emph{Complexity.} There are at most $(q+1)k^q$ charts. In one chart the quantified variables are the entries of $X$, $G$ and $x$, so

\[
N=2q\omega+2\omega\le2q^2+2q .
\]
The constraints are the rank-consistency disjunction over $t\le\min(\text{rows},\text{cols})$ of $M$ (at most $O(k)$ branches, each with $O(k)$ polynomials of degree $O(k)$), the $2r_0$ fixed-row equations, the free-row equations, and the Gram/spectral certificates (degree $2\omega+1$). Hence $s=O(k^2)$, $D=O(k)\le O(n)$, and all coefficients have bit length $\operatorname{poly}(n,\mathcal L)$ by the coefficient-size lemma of Appendix A.5. Renegar's Part II bound $B(\log B)(\log\log B)(sD)^{O(N)}$ gives

\[
(sD)^{O(N)}=\operatorname{poly}(n)^{O(q^2)}=\mathcal L^{O(q^2)}
\]
because $q$ is fixed; the chart factor $(q+1)k^q\le\mathcal L^{O(q)}$ is absorbed. The degenerate cases are listed in Section~\ref{sec:deg}.

The quantities entering the bound are collected here.

\begin{table*}[!t]
\centering
\caption{Quantities entering the fixed-order bound.}\label{tab:4}
\small
\begin{tabular}{p{0.450\textwidth}p{0.450\textwidth}}
\toprule
Quantity & Bound \\
\midrule
real variables $N$ & $2q\omega+2\omega\le2q^2+2q$ \\
pivot charts & $(q+1)k^q$ \\
rank-consistency branches $t$ & $O(k)$, one per possible rank \\
number of sign conditions $s$ & $O(k^2)$ \\
maximum degree $D$ & $O(k)$ \\
coefficient bit length $B$ & $\operatorname{poly}(n,\mathcal L)$ \\
construction size & $k^{O(q^2)}\operatorname{poly}(\mathcal L)$ \\
\bottomrule
\end{tabular}
\end{table*}
For fixed $q$ the exponent $O(q^2)$ is a constant, so the construction size is $\mathcal L^{O(q^2)}$ and Renegar's bound gives the same order. The rank index $t$ is a discrete branch, not a degree: it is counted in $s$, not folded into $D$.

Three representation details are worth recording.

\begin{enumerate}
\item \emph{Strict inequalities.} A condition $f>0$ is written as $\exists u:\ fu^2=1$, which is equivalent over the reals. There are $O(1)$ strict inequalities per branch, so this adds $O(1)$ variables and changes no order.
\item \emph{Disjunctions.} The rank-consistency disjunction is processed branch by branch: one real-feasibility query per $t$, with the running times added. This is why $t$ appears in $s$ and not in $D$.
\item \emph{Elementary symmetric polynomials.} $e_j(G)$ is represented by the coefficients of $\det(zI+G)$, computed by fraction-free elimination on the $k\times k$ polynomial matrix $zI+G$ in $O(k^3)$ polynomial operations of degree $O(k)$. The expansion is never written out term by term.
\end{enumerate}
\subsection{Degenerate cases}\label{sec:deg}
\begin{enumerate}
\item \emph{$k=0$} (that is, $\ker[C;CA]=0$). Then $N$ has no columns, so $H=H_A=0$ and both $(5a)$ and $(5b)$ hold for every $L$; in particular for $L=L_0$, so $q_{\min}=r_0$.
\item \emph{$\omega=0$} with $k>0$. Then $H=0$, and $\operatorname{rank}(\lambda H-H_A)=0$ for all $\lambda$ forces $H_A=0$; together with $H=0$ this is exactly $(5a)\wedge(5b)$ for that $L$. The branch $H=H_A=0$ of the encoding is therefore neither stronger nor weaker than the condition.
\item \emph{$q<r_0$}: rejected by definition, since $L$ must contain $L_0$.
\item \emph{$r=0$}: Proposition 3c (Section~\ref{sec:zero}).
\end{enumerate}
\section{Theorem 2: fixed effective dimension}\label{sec:thm2}
\textbf{Theorem 2 (bit model).} For rational input and every $r\ge0$, $q1$ is decidable in

\[
f(\kappa_{\rm obs})\operatorname{poly}(\mathcal L)
\]
bit operations, where

\[
\kappa_{\rm obs}=r\operatorname{rank}[N\ AN]
\]
and $r$ is the number of free observer rows. For $r=0$ the bound sharpens to $\operatorname{poly}(\mathcal L)$.

\end{thm}
\emph{Proof.} The two cases are separated, because the free-row parametrization degenerates at $r=0$.

\emph{Case 1: $r=0$.} There are no free rows, so $L=L_0$ is fixed and $\kappa_{\rm obs}=0$. Proposition 3c in Section~\ref{sec:zero} decides $(5a)$ and $(5b)$ for this single $L$ in $\operatorname{poly}(n,\mathcal L)$ bit operations. This is exactly the parameter value $0$, so the FPT statement has no boundary case left open.

\emph{Case 2: $r\ge1$.} Use Lemma B to parametrize all free rows by $\kappa_{\rm obs}$ real variables $t$. For each $\omega=0,\dots,q$, form the matrices $H(t),H_A(t)$, and impose

\[
G_H=H^{\mathsf T}H,\qquad G_S=H^{\mathsf T}H+H_A^{\mathsf T}H_A,
\]
\[
e_\omega(G_H)>0,\qquad e_{\omega+1}(G_S)=0,
\]
together with the $2\omega$-variable certificate

\[
B^{\mathsf T}Bx=e_0
\]
of Lemma D, where $B$ is built from the maximal minors of $P(\lambda)=\lambda H(t)-H_A(t)$ and $B^{\mathsf T}B$ is computed from $K(\lambda,\mu)$ by Cauchy--Binet. The branch $\omega=0$ is the one of Section~\ref{sec:deg}: there $H=0$ and $(5b)$ is $H_A=0$; Lemma D is not applied at $\omega=0$.

If a feasible $L=[L_0;R]$ exists, let $t$ be the coordinates of its free rows. Then $H=LN$, $H_A=LAN$; the Gram conditions encode $(5a)$ (Lemma B and the positivity of elementary symmetric functions), and the certificate encodes $(5b)$ (Lemma D). Hence the formula is true.

Conversely, if the formula is true for some $t$, define $L$ by taking $L_0$ and, for each free row pair in $\mathcal T$, a preimage under $[N\ AN]$. Then $H=LN$, $H_A=LAN$. The two Gram conditions give $\operatorname{rank}H=\omega$ and $\operatorname{row}H_A\subseteq \operatorname{row}H$, hence $(5a)$ by Lemma A; the certificate gives $(5b)$ by Lemma D. Thus the formula is equivalent to feasibility.

The quantified variables are

\[
N=\kappa_{\rm obs}+2\omega=f(\kappa_{\rm obs})
\]
(the discrete choice of $\omega$ is a finite disjunction, not an existential real variable). The number of polynomials, their degree, and their coefficient bit length are all $f(\kappa_{\rm obs})\operatorname{poly}(\mathcal L)$.

Renegar's Part II bit bound has the form

\[
B(\log B)(\log\log B)(sD)^{O(N)}
\]
for integer-coefficient input, where $N,s,D,B$ are the variable count, the number of polynomials, the maximum degree, and the maximum coefficient bit length. Substituting the bounds above gives

\[
f(\kappa_{\rm obs})\operatorname{poly}(\mathcal L)
\]
bit operations. This makes the bit-model statement unconditional, subject to the coefficient-size estimate in Appendix A.5.

\begin{cor}[2a] The same algorithm uses $f(\kappa_{\rm obs})\operatorname{poly}(\mathcal L)$ arithmetic operations.

\subsection{The zero-free-row case: an exact polynomial-time pencil test}\label{sec:zero}
When $r=0$ there are no free rows: $L=L_0$ is fixed and $\kappa_{\rm obs}=0$. Only two conditions remain, and both are decidable without any parameter-dependent search. $(5a)$ is a constant-rank test on $H=L_0N$, $H_A=L_0AN$. $(5b)$ is the constancy of the rank of the degree-one pencil $\lambda X-G$.

We first isolate the pencil statement.

\textbf{Lemma 8.1 (degree-one pencil test).} Let $X,G\in\mathbb Q^{q\times\omega}$ with $q\ge\omega\ge1$, and put $Q(\lambda)=\lambda X-G$. The following are equivalent:

\begin{enumerate}
\item $\operatorname{rank}Q(\lambda)=\omega$ for every $\lambda\in\mathbb C$;
\item the maximal $\omega\times\omega$ minors of $Q$ have no common complex root;
\item the maximal minors do not all vanish and their gcd in $\mathbb Q[\lambda]$ is a nonzero constant;
\item there is $S(\lambda)\in\mathbb Q[\lambda]^{\omega\times q}$ with $S(\lambda)Q(\lambda)=I_\omega$ and $\deg S\le2\omega-2$;
\item there exist matrices $S_0,\dots,S_{2\omega-2}\in\mathbb Q^{\omega\times q}$ satisfying the rational linear system \begin{equation*}
\tag{R0}
   -S_0G=I_\omega,\qquad
   S_{j-1}X-S_jG=0\quad(1\le j\le2\omega-2),\qquad
   S_{2\omega-2}X=0 .
\end{equation*}
\end{enumerate}
\end{cor}
\emph{Proof.} (1)$\Leftrightarrow$(2): $Q$ is $q\times\omega$, so $\operatorname{rank}Q(\lambda)<\omega$ holds exactly when every maximal minor vanishes at $\lambda$. (2)$\Leftrightarrow$(3): the minors have rational coefficients and $\mathbb C$ is algebraically closed, so a common complex root exists exactly when the gcd is nonconstant: a nonconstant gcd has a root in $\mathbb C$, and the minimal polynomial of a common root divides the gcd. (3)$\Rightarrow$(4): this is the adjugate/B\'ezout construction of Appendix B. For each $\omega$-element row set $I$, let $Q_I=E_IQ$ and $f_I=\det Q_I$; then $\operatorname{adj}(Q_I)E_IQ=f_II_\omega$, and B\'ezout coefficients with $\sum_Ig_If_I=1$, $\deg g_I<\omega$, give $S=\sum_Ig_I\operatorname{adj}(Q_I)E_I$. Since $\deg\operatorname{adj}(Q_I)\le\omega-1$, $\deg S\le2\omega-2$. (4)$\Rightarrow$(5): write $S=\sum_{j=0}^{2\omega-2}S_j\lambda^j$ and compare coefficients of $\lambda^k$ in $S(\lambda)(\lambda X-G)=I_\omega$; the coefficient of $\lambda^0$ is $-S_0G$, of $\lambda^k$ ($1\le k\le2\omega-2$) is $S_{k-1}X-S_kG$, and of $\lambda^{2\omega-1}$ is $S_{2\omega-2}X$. (5)$\Rightarrow$(1): evaluating the polynomial identity $S(\lambda)Q(\lambda)=I_\omega$ at any $\lambda_0\in\mathbb C$ exhibits a left inverse of $Q(\lambda_0)$, so $\operatorname{rank}Q(\lambda_0)=\omega$. $\square$

\emph{Size and bit complexity of (R0).} Writing $d=2\omega-2$, the unknowns are the entries of $d+1=2\omega-1$ matrices of size $\omega\times q$, i.e.

\[
N_{\rm var}=(2\omega-1)q\omega=O(q\omega^2)
\]
scalar unknowns; each of the $d+2=2\omega$ matrix equations contributes $\omega^2$ scalar equations, so

\[
N_{\rm eq}=2\omega^3 .
\]
All coefficients are entries of $X$, $G$, or $0,\pm1$. Clearing the denominators multiplies the system by an integer of bit length $O(\mathcal L_0)$, where $\mathcal L_0$ is the input length of $(X,G)$, and exact fraction-free elimination on the resulting integer system of size $M=O(\omega^3+q\omega^2)$, with entries of bit length $O(\mathcal L_0)$, runs in $\operatorname{poly}(M,\mathcal L_0)$ bit operations: the intermediate minors have bit length $O(M(\mathcal L_0+\log M))$ by Hadamard's inequality. Hence consistency of (R0) is decided in $\operatorname{poly}(q,\omega,\mathcal L_0)$ bit operations.

\emph{Remark on the construction.} The adjugate construction enumerates the $\binom q\omega$ maximal minors, which is exponential in $q$; it is used only to prove the existence of a certificate of degree at most $2\omega-2$. The algorithm never forms this sum: it searches for the coefficients of $S$ directly in (R0), whose size is polynomial. This is the step that turns a certificate-size statement into a decision algorithm.

\emph{Remark on higher degree.} The same argument applies to pencils of degree $d$: if the entries of $Q$ have degree at most $d$, then every maximal minor has degree at most $d\omega$, B\'ezout gives multipliers with $\deg g_I<d\omega$, and the adjugate contributes at most $d(\omega-1)$. Hence

\[
\deg S\le 2d\omega-d-1,
\]
which is $2\omega-2$ for $d=1$. Degree bounds for polynomial-matrix left inverses are studied classically \cite{fab90}; the bound above is the elementary adjugate--B\'ezout one, and only the case $d=1$ is used in this paper. The general bound may be useful if the observer model is extended. (For $d=0$ the pencil is constant and the statement is the obvious rank test.)

\begin{prop}[3c] Let $r=0$, so $L=L_0$. Feasibility of $(5a)\wedge(5b)$ is decidable in $\operatorname{poly}(n,\mathcal L)$ bit operations.

\emph{Algorithm.} Compute $N$ (a basis of $\ker[C;CA]$), $H=L_0N$, $H_A=L_0AN$.

\begin{enumerate}
\item Check $\operatorname{rank}[H_A;H]=\operatorname{rank}H$; if it fails, reject, since this is $(5a)$ (Lemma A, Appendix A.1).
\item Put $\omega=\operatorname{rank}H$. If $\omega=0$, then $H=0$ and $(5b)$ reads $\operatorname{rank}(\lambda H-H_A)=0$ for all $\lambda$, i.e. $H_A=0$; accept exactly in that case.
\item Otherwise take a rank factorization $H=XB$ with $X\in\mathbb Q^{q\times\omega}$ of full column rank and $B\in\mathbb Q^{\omega\times k}$ of full row rank. Step 1 gives $\operatorname{row}H_A\subseteq\operatorname{row}H$ (Lemma A), hence $H_A=FH$ for some $F$; set $G=FX$, so $H_A=GB$.
\item Run the test of Lemma 8.1 on $(X,G)$.
\end{enumerate}
\emph{Correctness.} Since $B$ is of full row rank, $\operatorname{rank}(\lambda H-H_A)=\operatorname{rank}((\lambda X-G)B) =\operatorname{rank}(\lambda X-G)$ for every $\lambda$, and $(5b)$ is $\operatorname{rank}(\lambda H-H_A)=\omega$ for all $\lambda$. Lemma 8.1 therefore decides $(5b)$. Step 1 is $(5a)$, and step 2 is the exact $\omega=0$ reading of the same identity. $\square$

\emph{Complexity.} Step 1 is one rank computation; step 3 is a rank factorization and an exact linear solve; step 4 is $\operatorname{poly}(q,\omega,\mathcal L_0)\subseteq\operatorname{poly}(n,\mathcal L)$ by Lemma 8.1. All steps use rational arithmetic on data of polynomial bit length.

\emph{Interpretation in Kronecker canonical form.} The pencil $Q$ has full column rank at every finite complex value if and only if its Kronecker form (in the sense of \cite{bee88}) contains no right singular block and no finite regular elementary-divisor block. Left singular blocks and \emph{infinite} regular elementary divisors are permitted: they do not produce a rank drop at any finite $\lambda$. For instance $Q(\lambda)=1$ has an infinite elementary divisor of size one and satisfies the condition; so does $Q(\lambda)=I_2$, and so does the nonconstant pencil $Q(\lambda)=\begin{bmatrix}1&\lambda\\0&1\end{bmatrix}$, whose determinant is the constant $1$. Equivalently, in the language of Appendix B, the condition is that the ideal generated by the maximal minors is the unit ideal; only the \emph{finite} spectrum is constrained.

\emph{Witness.} If (R0) is consistent, exact elimination returns $S_0,\dots,S_{2\omega-2}$, hence an explicit polynomial left inverse $S(\lambda)$, verifiable by multiplying the two polynomial matrices. The test is therefore not merely a yes/no decision.

\emph{Lean status.} The left-inverse/injectivity core of this section is verified in \texttt{Q1Paper.FPT} (\texttt{injective\_\allowbreak{}of\_\allowbreak{}leftInverse}, \texttt{finrank\_\allowbreak{}le\_\allowbreak{}of\_\allowbreak{}leftInverse}, \texttt{kappaObs\_\allowbreak{}zero}). The linear system (R0), its counts, and the bit-complexity statement are paper-level; see Appendix C for the exact checks.

\section{Theorem 3 and the combined bound}\label{sec:thm3}
\textbf{Theorem 3.} For every fixed $q$, $\mathrm{DEC}(q)\in\exists\mathbb R$. Consequently the threshold problem $\mathrm{MINDEC}(Q)$ is in $\exists\mathbb R$ for every $Q$.

\end{prop}
\emph{Proof.} By Appendix B: the encoding is a finite disjunction over $\omega=0,\dots,\min(q,k)$, each branch a polynomial-size system with the data constraints (B.1)--(B.2), and its solutions correspond exactly to feasible $L$. For $\mathrm{MINDEC}(Q)$ take the finite disjunction over $q=r_0,\dots,Q$. $\square$

\begin{cor}[4] For every $r\ge0$, the decision time in the bit model satisfies

\[
T\le\min\{\mathcal L^{O(q^2)},\ f(\kappa_{\rm obs})\operatorname{poly}(\mathcal L)\}.
\]
For $r=0$ the second bound is $\operatorname{poly}(\mathcal L)$ by Proposition 3c.

\end{cor}
\emph{Proof.} Theorem 1 gives the first bound for fixed $q$; Theorem 2 the second. Both algorithms decide the same language, so one may run either. $\square$

\section{Discussion and open problems}\label{sec:disc}
The two upper bounds are parameterized: the exponent in Theorem 1 depends on $q$, and the factor in Theorem 2 depends on $\kappa_{\rm obs}$. Neither parameter is bounded in the unrestricted problem, so the results do not imply $q1\in\mathrm P$.

Three outcomes remain compatible with the present results: $q1\in\mathrm P$; $q1$ is $\exists\mathbb R$-complete; or $q1$ has an intermediate parameterized classification. The paper makes no claim among them.

\subsection{Limitations}\label{sec:limits}
\begin{enumerate}
\item \textbf{The $r=0$ branch: resolved, with a residual audit.} For $r=0$ the two conditions reduce to a constant-rank test and the constancy of the rank of the degree-one pencil $\lambda X-G$. Section~\ref{sec:zero} decides the latter by a finite rational linear system derived from the degree bound $2\omega-2$; the test is deterministic and exact, with $\operatorname{poly}(q,\omega)$ unknowns and equations, and it returns an explicit witness. The remaining issue is to complete the explicit bit-complexity audit: the constants in the bit bound of Section~\ref{sec:zero} are standard Hadamard/elimination bounds but have not been checked line by line.
\item \textbf{Coefficient bounds.} The coefficient-size lemma in Appendix A.5 is paper-level; its constants have not been independently audited.
\item \textbf{Hardness.} The paper contains no hardness reduction. The unrestricted complexity of $q1$ remains open.
\end{enumerate}
\subsection{The size of the parameter}\label{sec:paramsize}
$\kappa_{\rm obs}=r\operatorname{rank}[N\ AN]\le 2rn$ can be quadratic in $n$; when it is, the bound of Theorem 2 is not useful in practice. The theorem is therefore best read as a statement about the \emph{reduction}: $q1$ reduces, in time $f(\kappa_{\rm obs})\operatorname{poly}(\mathcal L)$, to a real-feasibility problem whose dimension is the number $rd_0$ of free real coordinates. Two remarks.

\begin{enumerate}
\item \emph{The reduction is where the work is.} Building the formula needs no enumeration of the $\binom q\omega\binom k\omega$ maximal minors (Cauchy--Binet, Section~\ref{sec:cb}), no enumeration of the $r_0$ fixed rows, and no chart over $\operatorname{Gr}(\omega,k)$; the construction cost is $f(\kappa_{\rm obs})\operatorname{poly}(\mathcal L)$, which is what separates the result from the trivial statement "$q1$ is a $\kappa_{\rm obs}$-variable real problem".
\item \emph{The parameter is bounded on infinite families.} Take $p=n-1$, $C=[I_{n-1}\ 0]$, and $A$ whose last column lies in $\ker C$ (so that the last column of $CA$ vanishes). Then $\ker[C;CA]=\operatorname{span}\{e_n\}$, hence $k=1$, and $d_0=\operatorname{rank}[N\ AN]\le2$; with one free row, $\kappa_{\rm obs}\le2$ while $n$ is unbounded. On this family Theorem 2 gives $\operatorname{poly}(\mathcal L)$ time with no $n$-dependence in the exponent (checked for $n=4,\dots,8$ by \texttt{c319}).
\end{enumerate}
Whether $\kappa_{\rm obs}$ is optimal is open, and three levels of further reduction are worth separating.

\begin{enumerate}
\item \emph{The kernel quotient.} $d_0=\dim\operatorname{im}\Phi$ with $\Phi(v)=(vN,vAN)$ is already the \emph{exact} dimension of the free-pair space, so no further reduction is possible at this level.
\item \emph{Observer-row transformations.} The conditions depend on row spaces, hence are invariant under row operations of $L$. Those preserving $L_0$ are block lower-triangular, $T=\begin{bmatrix}I_{r_0}&0\\ U&S\end{bmatrix}$ with $\det S\ne0$; the free coordinates may therefore contain redundant presentations of the same observer subspace. A quotient by this symmetry could lower the parameter.
\item \emph{Spectral invisibility.} One may hope to divide by directions that change neither the rank condition nor the finite spectrum. Such a quotient must preserve the interaction between $H$ and $H_A$, not only $\operatorname{rank}H$; a quotient preserving $(5a)$ but changing the spectrum would be invalid.
\end{enumerate}
A natural first refinement is not a spectral quotient but an affine-dimension count: some free-row directions may already be constrained by the fixed rows $L_0$, and the effective rank of the free-row map after imposing the fixed-row equations could be strictly smaller than $rd_0$. This is the first step to take; the existence of a spectral quotient, and the optimality of $\kappa_{\rm obs}$, are open.

Growth of the parameter-dependent factor. Renegar's bound is $(sD)^{O(N)}$ with $sD=\operatorname{poly}(\kappa_{\rm obs})$ and $N=\kappa_{\rm obs}+2\omega$, so $f(\kappa_{\rm obs})$ is roughly $\kappa_{\rm obs}^{O(\kappa_{\rm obs})}$. The result is primarily structural: it identifies a parameter governing the complexity bound, although the dependence on that parameter may limit practical applicability for large parameter values.

\appendices
\section{Proof details}\label{app:a}
\subsection*{A.1 Matrix-to-subspace bridge for Lemma A}
Let $K_0=\ker[C;CA]$, let $N$ be a basis of $K_0$, and let $H=LN$, $H_A=LAN$. Write $\mathcal J=\operatorname{row}[C;CA]$. The annihilator identity

\[
K_0^{\perp}=\mathcal J
\]
holds for the standard bilinear form: a row vector vanishes on $K_0$ exactly when it is a linear combination of the rows of $[C;CA]$.

The rank equality $(5a)$ is equivalent to $\operatorname{row}(LA)\subseteq\operatorname{row}[L;CA;C]$. To pass to the compressed matrices, use the identity

\[
\operatorname{rank}\begin{bmatrix}J\\ M\end{bmatrix}
=\operatorname{rank}J+\operatorname{rank}(MN),
\qquad N\text{ a basis of }\ker J,
\]
which holds because restriction of row functionals to $\ker J$ has kernel $\operatorname{row}J$. Apply it with $J=[C;CA]$ and with $M=[LA;L]$. Then

\[
\operatorname{rank}[LA;L;CA;C]
=\operatorname{rank}J+\operatorname{rank}[H_A;H],
\]
while

\[
\operatorname{rank}[L;CA;C]
=\operatorname{rank}J+\operatorname{rank}H.
\]
Hence $(5a)$ is equivalent to $\operatorname{rank}[H_A;H]=\operatorname{rank}H$, i.e. to $\operatorname{row}H_A\subseteq\operatorname{row}H$.

Conversely, assume $H_A=FH$. Then

\[
(LA-FL)N=0.
\]
Hence every row of $LA-FL$ annihilates $K_0$, so every such row lies in $K_0^{\perp}=\mathcal J\subseteq\operatorname{row}[L;CA;C]$. Therefore $\operatorname{row}(LA)\subseteq\operatorname{row}[L;CA;C]$, which is $(5a)$.

For $(5b)$, apply the same identity to the pencil $M(\lambda)=[\lambda L-LA;CA;C]$ and $R=[L;CA;C]$. Since $N$ is a basis of the common kernel of $CA,C$,

\[
\operatorname{rank}M(\lambda)
=\operatorname{rank}J+\operatorname{rank}(\lambda H-H_A),
\qquad
\operatorname{rank}R=\operatorname{rank}J+\operatorname{rank}H.
\]
Equivalently,

\[
\operatorname{rank}M(\lambda)=\operatorname{rank}R
\iff
\operatorname{rank}(\lambda H-H_A)=\operatorname{rank}H.
\]
Thus $(5b)$ is equivalent to the pencil $P(\lambda)=\lambda H-H_A$ having constant rank $\omega=\operatorname{rank}H$. A rectangular pencil has constant rank $\omega$ on $\mathbb C$ exactly when its $\omega\times\omega$ maximal minors have no common root, i.e. when their gcd is a nonzero constant. This proves Lemma A(iii); Lemma A(iv) follows at $\lambda=0$.

\subsection*{A.2 Lemma B: free-row parametrization}
For a free row $r_i\in\mathbb R^{1\times n}$,

\[
(H_i,(H_A)_i)=r_i[N\ AN].
\]
Thus every free row pair lies in $\mathcal T=\operatorname{row}[N\ AN]$. The linear map $r\mapsto r[N\ AN]$ has image $\mathcal T$, so its rank is $d_0=\dim\mathcal T$. Choosing a basis of $\mathcal T$, each free row is described by exactly $d_0$ coordinates, and $r$ free rows by $rd_0=\kappa_{\rm obs}$ coordinates.

Because $N$ has full column rank, its $k$ columns are linearly independent in the stacked matrix $[N\ AN]$; hence $k\le d_0$. For $r\ge1$, $d_0\le rd_0=\kappa_{\rm obs}$.

\subsection*{A.3 Lemma C: rank factorization and consistency}
Let $\omega>0$ and choose a pivot set $J$, $|J|=\omega$, such that $X=H_{:,J}$ has full column rank. Then $H=XB$ for a unique $B\in\mathbb C^{\omega\times k}$ with $B_{:,J}=I_\omega$: the columns in $J$ determine the coordinates, and the remaining columns are the coordinates of the corresponding columns of $H$ in the basis $X$.

If $H_A=FH$, set $G=FX$. Then $H_A=GB$. Conversely, if $H_A=GB$ and $X_LX=I_\omega$, then $F=GX_L$ satisfies $FH=GX_LXB=GB=H_A$. This is the algebraic content of the substitution $G=FX$; it is Lean-verified in \texttt{Q1Paper.G\_\allowbreak{}form\_\allowbreak{}iff}.

Because $B$ has full row rank (its columns $J$ form $I_\omega$),

\[
\operatorname{rank}P(\lambda)
=\operatorname{rank}((\lambda X-G)B)
=\operatorname{rank}(\lambda X-G).
\]
For the linear system $Mb=c$, consistency is equivalent to $c\in\operatorname{im}M$. Over the reals,

\[
\operatorname{im}M=\operatorname{im}(MM^{\mathsf T}),
\qquad
\operatorname{im}[M\ c]=\operatorname{im}(MM^{\mathsf T}+cc^{\mathsf T}).
\]
Hence

\[
\exists b:\ Mb=c
\iff
\operatorname{rank}(MM^{\mathsf T})
=\operatorname{rank}(MM^{\mathsf T}+cc^{\mathsf T}).
\]
For a positive semidefinite matrix $\Gamma$, the $t$-th elementary symmetric polynomial satisfies $e_t(\Gamma)>0$ and $e_{t+1}(\Gamma)=0$ exactly when $\operatorname{rank}\Gamma=t$. Applying this to $\Gamma=MM^{\mathsf T}$ and $\Gamma'=MM^{\mathsf T}+cc^{\mathsf T}$, and using $\Gamma'\succeq\Gamma$, gives the displayed disjunction in Lemma C.

\subsection*{A.4 Lemma D: Bezout, Gram certificate, Cauchy--Binet}
\textbf{Bezout.} Let $f_1,\dots,f_s\in\mathbb R[\lambda]$ have degree at most $\omega$ and no common complex root. Choose a nonzero real combination $g=\sum a_i f_i$. If $g$ is constant, then $g^{-1}a_i$ are constant multipliers. Otherwise $g$ has finitely many complex roots. At every root $\alpha$ of $g$, not all $f_i(\alpha)$ vanish, so the condition $h(\alpha)=0$ for $h=\sum b_i f_i$ is a proper real linear condition on $b$. Avoiding the finite union of these proper subspaces gives $h=\sum b_i f_i$ with no common root with $g$. Then $g,h$ are coprime. The Euclidean Bezout identity gives $ug+vh=1$ with $\deg u<\deg h\le\omega$ and $\deg v<\deg g\le\omega$. Substituting the definitions of $g,h$,

\[
1=\sum_i (u a_i+v b_i) f_i,
\]
and each multiplier $u a_i+v b_i$ has degree $<\omega$, because $a_i,b_i$ are constants and $\deg u,\deg v<\omega$. This is the two-polynomial core verified in \texttt{Q1Paper.bezout\_\allowbreak{}of\_\allowbreak{}gcd\_\allowbreak{}eq\_\allowbreak{}one}; the degree bookkeeping is paper-level.

\textbf{Gram certificate.} Let $f_i$ be the maximal minors of $P(\lambda)$, indexed by $i=(I,J)$; each has degree at most $\omega$. Let $B$ have rows indexed by $(i,j)$, $0\le j<\omega$, columns indexed by $t=0,\dots,2\omega-1$, and entries

\[
B_{(i,j),t}=c_{i,t-j},
\]
where $c_{i,a}$ is the coefficient of $\lambda^a$ in $f_i$ and $c_{i,a}=0$ outside $0\le a\le\omega$. Thus the number of rows is $\omega\binom{k}{\omega}$ and the number of columns is $2\omega$. The row indexed by $(i,j)$ is the coefficient vector of $\lambda^j f_i$, whose degree is at most $j+\deg f_i\le2\omega-1$. Consequently

\[
\operatorname{row}B
=\left\{\text{coefficient vector of }\sum_i g_i f_i:
\deg g_i<\omega\right\}.
\]
Therefore

\[
e_0\in\operatorname{row}B
\iff
\exists g_i,\ \deg g_i<\omega,\ \sum_i g_i f_i=1.
\]
By the Bezout lemma, the right-hand side is equivalent to $(5b)$. For real $B$, $\operatorname{row}B=\operatorname{im}(B^{\mathsf T}B)$, so this is further equivalent to the linear system $B^{\mathsf T}Bx=e_0$.

\textbf{Cauchy--Binet.} We use the rectangular identity

\[
\det(AC)=\sum_{|J|=\omega}\det(A_{:,J})\det(C_{J,:}),
\qquad A\in\mathbb R^{\omega\times k},\ C\in\mathbb R^{k\times\omega}.
\]
For each column set $J$, apply it to $A=P(\lambda)_{:,J}^{\mathsf T}$ and $C=P(\mu)_{:,J}$; the sum over row sets $I$ gives

\[
\det\left(P(\lambda)_{:,J}^{\mathsf T}P(\mu)_{:,J}\right)
=\sum_{|I|=\omega}
f_{I,J}(\lambda)f_{I,J}(\mu).
\]
Summing over $J$,

\[
K(\lambda,\mu)=\sum_{|J|=\omega}
\det\left(P(\lambda)_{:,J}^{\mathsf T}P(\mu)_{:,J}\right)
=\sum_{I,J}f_{I,J}(\lambda)f_{I,J}(\mu)
=\sum_i f_i(\lambda)f_i(\mu).
\]
If $S_{a,b}$ is the coefficient of $\lambda^a\mu^b$ in $K$, then

\[
(B^{\mathsf T}B)_{t,t'}=\sum_{j=0}^{\omega-1}S_{t-j,t'-j}.
\]
Only $\binom{k}{\omega}\le2^{\kappa_{\rm obs}}$ minors of order $\omega$ occur. This is the content of \texttt{c313}; the general Cauchy--Binet identity is paper-level because Mathlib currently has no rectangular Cauchy--Binet theorem.

\subsection*{A.5 Complexity estimates}
We use the standard bit model: the input has rational entries of total bit length $\mathcal L$, and $n,p,q\le n\le\mathcal L$. Renegar's Part II bit bound \cite{ren92} has the form

\[
B(\log B)(\log\log B)(sD)^{O(N)},
\]
where $N$ is the number of real variables, $s$ the number of polynomials, $D$ the maximum degree, and $B$ the maximum coefficient bit length. The exponential factor depends only on $N$; the coefficient-length contribution is polynomial in $\log B$.

\textbf{Coefficient-size lemma.} Let the input entries have bit length at most $\mathcal L$, and let $n\le\mathcal L$. Then:

\begin{enumerate}
\item the kernel basis $N$, the matrix $D$, and the entries of $H,H_A,G_H,G_S$ have bit length $\operatorname{poly}(n,\mathcal L)$;
\item every maximal minor $f_i$ of $P(\lambda)=\lambda H-H_A$, and every coefficient of $f_i$, has bit length $\operatorname{poly}(n,\mathcal L)$;
\item the entries of $B^{\mathsf T}B$ and of $e_j(G_H),e_j(G_S)$ for $j\le\kappa_{\rm obs}+1$ have bit length $f(\kappa_{\rm obs})\operatorname{poly}(\mathcal L)$.
\end{enumerate}
\emph{Proof.} Gaussian elimination on $[C;CA]$ writes $N$ as ratios of minors of the input; Cramer's rule and Hadamard's bound give $\operatorname{poly}(n,\mathcal L)$ bits. The same applies to $D$. Each entry of $H=LN$ is a sum of $n$ products, and each Gram entry is a sum of $q$ such products, so the bit length remains $\operatorname{poly}(n,\mathcal L)$. A maximal minor of $P$ is a determinant of an $\omega\times\omega$ matrix whose entries are linear in the entries of $H,H_A$; Hadamard's bound gives $\operatorname{poly}(n,\mathcal L)$ bits for each coefficient. The matrix $B^{\mathsf T}B$ has at most $2^{\kappa_{\rm obs}}$ determinant terms, and $e_j$ has at most $k^{j}\le\kappa_{\rm obs}^{\kappa_{\rm obs}}$ terms; summing gives $f(\kappa_{\rm obs})\operatorname{poly}(\mathcal L)$. $\square$

\subsubsection*{Fixed $q$}
For fixed $q$, the algorithm enumerates at most $(q+1)k^q$ pivot charts. In one chart the quantified variables are the entries of $X,G$ and the spectral certificate, so

\[
N=q\omega+q\omega+2\omega=2q\omega+2\omega\le2q^2+2q .
\]
The constraints are the rank-consistency disjunction (at most $O(k)$ branches) and the Gram/spectral certificates; hence $s=O(k^2)$, $D=O(k)\le O(n)$, and the coefficients have bit length $\operatorname{poly}(n,\mathcal L)$ by the coefficient-size lemma above. Renegar's bound $B(\log B)(\log\log B)(sD)^{O(N)}$ therefore gives

\[
(sD)^{O(N)}=\operatorname{poly}(n)^{O(q^2)}=\mathcal L^{O(q^2)},
\]
since $q$ is fixed, and the chart enumeration contributes the factor $(q+1)k^q\le\mathcal L^{O(q)}$. Hence

\[
\mathcal L^{O(q^2)}
\]
bit operations. The cases $\omega=0$, $k=0$, $q<r_0$ and $r=0$ are treated in Section~\ref{sec:deg} (the last by Proposition 3c).

\subsubsection*{Fixed $\kappa_{\rm obs}$}
Assume $r\ge1$, so $k\le d_0\le\kappa_{\rm obs}$; both inequalities are Lean-verified (\texttt{finrank\_\allowbreak{}domain\_\allowbreak{}le\_\allowbreak{}freePairSpace}, \texttt{freePairDim\_\allowbreak{}le\_\allowbreak{}kappaObs}). The Gram matrices $G_H=H^{\mathsf T}H$ and $G_S=H^{\mathsf T}H+H_A^{\mathsf T}H_A$ are $k\times k$, hence of size at most $\kappa_{\rm obs}$. Their elementary symmetric polynomials $e_j$ for $j\le\omega+1\le\kappa_{\rm obs}+1$ can be evaluated by determinant expansion in $k^{O(j)}=f(\kappa_{\rm obs})$ operations.

The matrix $B^{\mathsf T}B$ is computed from the Cauchy--Binet kernel $K(\lambda,\mu)$. The sum has at most $\binom{k}{\omega}\le2^{\kappa_{\rm obs}}$ terms, each an $\omega\times \omega$ determinant; all coefficients are rational of bit length $f(\kappa_{\rm obs})\operatorname{poly}(\mathcal L)$. Thus the preprocessing costs $f(\kappa_{\rm obs})\operatorname{poly}(\mathcal L)$.

The real-feasibility problem has

\[
N=\kappa_{\rm obs}+2\omega=f(\kappa_{\rm obs})
\]
variables, at most $2\omega+O(1)=f(\kappa_{\rm obs})$ equations, degree at most $\kappa_{\rm obs}+1$, and coefficient bit length $f(\kappa_{\rm obs})\operatorname{poly}(\mathcal L)$. Renegar's Part II bound $B(\log B)(\log\log B)(sD)^{O(N)}$ gives

\[
f(\kappa_{\rm obs})\operatorname{poly}(\mathcal L)
\]
bit operations. The case $r=0$ is not covered by this counting; it is decided separately by Proposition 3c, whose own counting is in Section~\ref{sec:zero}. The cases $k=0$ and $\omega=0$ are direct.

\subsubsection*{A.5.1 The $r=0$ branch: proof details}
Lemma 8.1 rests on the adjugate construction for \emph{rectangular} pencils, so we spell out why no square assumption enters. For $Q\in\mathbb Q[\lambda]^{q\times\omega}$ and each $\omega$-element row set $I\subseteq\{1,\dots,q\}$, let $E_I\in\mathbb Q^{\omega\times q}$ be the row selection and $Q_I=E_IQ$. Then

\[
\operatorname{adj}(Q_I)E_IQ=f_II_\omega,\qquad f_I=\det Q_I,
\]
which is the identity used in Appendix B; hence a B\'ezout relation $\sum_Ig_If_I=1$ with $\deg g_I<\omega$ yields the polynomial left inverse $S=\sum_Ig_I\operatorname{adj}(Q_I)E_I$. Since $\deg\operatorname{adj}(Q_I)\le\omega-1$, the degree bound $\deg S\le2\omega-2$ holds for $q\ge\omega$, with no square case assumed. Two points are worth recording.

\begin{enumerate}
\item \emph{Existence versus algorithm.} The sum defining $S$ runs over $\binom q\omega$ row sets and is exponentially large; it is used only to prove that a certificate of degree at most $2\omega-2$ exists. The algorithm solves the linear system (R0) instead, which has $O(q\omega^2)$ unknowns and $2\omega^3$ equations.
\item \emph{$(5a)$ comes first.} Proposition 3c needs $\operatorname{row}H_A\subseteq\operatorname{row}H$ in order to produce $G$ with $H_A=GB$. If $(5a)$ fails, the instance is infeasible and the spectral question is not asked; the zero-free-row decision is therefore the conjunction of one constant-rank test and one pencil test.
\end{enumerate}
The Kronecker form is not used by the algorithm. The reading recorded in Section~\ref{sec:zero} only corrects an earlier statement: infinite regular elementary divisors are compatible with the condition, because they do not produce a rank drop at a finite $\lambda$.

\section{Real-algebraic certificate}\label{app:b}
For fixed $q$, $(5b)$ is the condition that the $q\times\omega$ pencil

\[
Q(\lambda)=\lambda X-G
\]
has rank $\omega$ for every $\lambda\in\mathbb C$. This is stronger than having rank $\omega$ over the fraction field $\mathbb C(\lambda)$: for example $Q(\lambda)=\binom{\lambda}{\lambda}$ has fraction-field rank one but drops rank at $\lambda=0$.

By the univariate Quillen--Suslin/Bezout argument, $Q$ has rank $\omega$ at every $\lambda$ iff its maximal minors generate the unit ideal in $\mathbb C[\lambda]$, equivalently iff there is a polynomial left inverse

\[
S(\lambda)Q(\lambda)=I_\omega.
\]
The degree bound is $2\omega-2$. For each $\omega$-element row set $I$, let $Q_I$ be the corresponding $\omega\times\omega$ submatrix and $f_I=\det Q_I$. If $Q$ has rank $\omega$ at every $\lambda$, the $f_I$ have gcd one. Lemma D1 gives $g_I$ with $\deg g_I<\omega$ and $\sum_I g_I f_I=1$. Let $S_I$ be the $\omega\times q$ matrix whose $I$-columns are $\operatorname{adj}(Q_I)$ and whose other columns are zero. Then

\[
S_I Q=f_I I_\omega,
\]
so $S=\sum_I g_I S_I$ satisfies $SQ=I_\omega$. Since $\deg\operatorname{adj}(Q_I)\le\omega-1$ and $\deg g_I<\omega$, $\deg S\le2\omega-2$.

For the $\exists\mathbb R$ encoding we do not enumerate the minors. For each $\omega=0,\dots,\min(q,k)$, introduce

\[
X\in\mathbb R^{q\times\omega},\quad
B\in\mathbb R^{\omega\times k},\quad
G\in\mathbb R^{q\times\omega},\quad
U\in\mathbb R^{\omega\times q},\quad
V\in\mathbb R^{k\times\omega},
\]
and the coefficients of $S(\lambda)=\sum_{j=0}^{2\omega-2}S_j\lambda^j$, where $S_j\in\mathbb R^{\omega\times q}$. Impose

\[
H=XB,\qquad H_A=GB,\qquad UX=I_\omega,\qquad BV=I_\omega,
\]
and the polynomial identity

\[
S(\lambda)(\lambda X-G)=I_\omega,
\]
which is linear in the coefficients $S_j$. The rank conditions $UX=I_\omega$ and $BV=I_\omega$ are polynomial equations; $BV=I_\omega$ makes $B$ full row rank, hence

\[
\operatorname{rank}P(\lambda)=\operatorname{rank}((\lambda X-G)B)
=\operatorname{rank}(\lambda X-G)=\omega
\]
for all $\lambda$. Also $H_A=GB=(GU)XB=FH$ with $F=GU$, so $(5a)$ holds.

It remains to tie $(H,H_A)$ to the input. The rows of $H=LN$, $H_A=LAN$ satisfy, with $h_i^0=(L_0N)_i$ and $g_i^0=(L_0AN)_i$,

\begin{equation*}
X_iB=h_i^0,\qquad G_iB=g_i^0\qquad(1\le i\le r_0),
\tag{B.1}
\end{equation*}
and, for the free rows, $(H_i,H_{A,i})\in\mathcal T=\operatorname{row}[N\ AN]$, equivalently

\begin{equation*}
D\begin{bmatrix}(X_iB)^{\mathsf T}\\ (G_iB)^{\mathsf T}\end{bmatrix}=0
\qquad(r_0<i\le q),
\tag{B.2}
\end{equation*}
where $D$ is a rational matrix with $\ker D=\mathcal T$. Both are polynomial (in fact bilinear) equations. They are not optional: without (B.1)--(B.2) the encoding would accept arbitrary pairs $(H,H_A)$.

The correctness splits into two independent statements.

\emph{Soundness.} Suppose the system has a solution. Set $L$ with first $r_0$ rows $L_0$ and, for each $i>r_0$, a preimage $r_i$ with $(X_iB,G_iB)=r_i[N\ AN]$, which exists by (B.2). Then $H=LN$, $H_A=LAN$ by (B.1)--(B.2); $(5a)$ follows from $H_A=GB=(GU)XB$; and $(5b)$ from $\operatorname{rank}(\lambda H-H_A)=\operatorname{rank}(\lambda X-G)=\omega$ for all $\lambda$, which uses $BV=I_\omega$ (full row rank of $B$) and $S(\lambda)(\lambda X-G)=I_\omega$.

\emph{Completeness.} A feasible $L$ supplies $X,B,U,V$ from a rank factorization $H=XB$ with $B_{:,J}=I_\omega$, $G=FX$ with $H_A=FH$ (Lemma A), and $S$ from the left-inverse construction of Appendix B; then (B.1)--(B.2) hold by construction and the remaining equations are identities.

The two constraints $UX=I_\omega$ and $BV=I_\omega$ are not redundant bookkeeping: the first produces the factorization $H_A=FH$, the second makes $\operatorname{rank}((\lambda X-G)B)=\operatorname{rank}(\lambda X-G)$ hold at every $\lambda$.

The $\omega=0$ case is handled by $H=H_A=0$, which (B.1) forces when $L_0N=0$ and rejects otherwise; it is a separate algebraic branch, not a degree-negative instance of (R0). The number of variables is

\[
O(q\omega+\omega k+q\omega+\omega q+k\omega+(2\omega-1)\omega q)
=O(n^3),
\]
the number of equations is polynomial, and the degrees are bounded by $2\omega-2\le2n$. The disjunction over $\omega$ has at most $n+1$ branches. Hence the encoding has polynomial size. This proves $\mathrm{DEC}(q)\in\exists\mathbb R$ for every fixed $q$, and hence $\mathrm{MINDEC}(Q)\in\exists\mathbb R$.

\section{Verification record}\label{app:c}
See \texttt{paper/verification\_\allowbreak{}log.md}.

\textbf{Reproducibility.} The Lean sources are mirrored under \texttt{paper/lean/Q1Paper/}; the WSL build command is \texttt{bash scripts/check.sh}. The exact Python checks are under \texttt{tools/verification/} and are summarized in \texttt{paper/verification\_\allowbreak{}log.md}. These checks are finite evidence; they are not proofs and do not replace the paper arguments or the Lean kernel.

Lean declarations currently covering the algebraic kernel:

\begin{itemize}
\item \texttt{Q1Paper.ker\_\allowbreak{}compressed\_\allowbreak{}iff} (Lemma A compression criterion);
\item \texttt{Q1Paper.mem\_\allowbreak{}freePairSpace\_\allowbreak{}iff}, \texttt{finrank\_\allowbreak{}freePairSpace\_\allowbreak{}add\_\allowbreak{}ker}, \texttt{kappaObs\_\allowbreak{}eq\_\allowbreak{}mul\_\allowbreak{}freePairSpace} (Lemma B free-pair space);
\item \texttt{Q1Paper.exists\_\allowbreak{}solution\_\allowbreak{}iff\_\allowbreak{}span\_\allowbreak{}le\_\allowbreak{}range}, \texttt{exists\_\allowbreak{}solution\_\allowbreak{}iff\_\allowbreak{}range\_\allowbreak{}eq\_\allowbreak{}sup\_\allowbreak{}span} (Lemma C consistency);
\item \texttt{Q1Paper.bezout\_\allowbreak{}of\_\allowbreak{}gcd\_\allowbreak{}eq\_\allowbreak{}one}, \texttt{Q1Paper.sum\_\allowbreak{}sq\_\allowbreak{}eq\_\allowbreak{}zero\_\allowbreak{}iff} (Lemma D Bezout/positivity core);
\item \texttt{Q1Paper.G\_\allowbreak{}form\_\allowbreak{}iff} (Lemma C \texttt{G=FX} substitution);
\item \texttt{Q1Paper.injective\_\allowbreak{}of\_\allowbreak{}leftInverse}, \texttt{finrank\_\allowbreak{}le\_\allowbreak{}of\_\allowbreak{}leftInverse}, \texttt{kappaObs\_\allowbreak{}zero} (Lemma F r=0 core);
\item \texttt{Q1Paper.pairMap\_\allowbreak{}apply}, \texttt{Q1Paper.ker\_\allowbreak{}pairMap}, \texttt{Q1Paper.finrank\_\allowbreak{}range\_\allowbreak{}pairMap}, \texttt{Q1Paper.finrank\_\allowbreak{}range\_\allowbreak{}pairMap\_\allowbreak{}common};
\item \texttt{b24\_\allowbreak{}factor}, \texttt{b24\_\allowbreak{}factor\_\allowbreak{}unique} (reused B.24 quotient-factor lemmas, rebuilt in the new project);
\item \texttt{KappaObs.kappaObs\_\allowbreak{}eq\_\allowbreak{}row\_\allowbreak{}count} (reused effective-dimension lemma).
\end{itemize}
The complexity/QE steps are not Lean-formalized; see \texttt{LEAN\_\allowbreak{}MAP.md}.

The r=0 rectangular-adjugate construction is checked on 300 random full-row-rank linear pencils by \texttt{tools/verification/c315\_\allowbreak{}r0\_\allowbreak{}left\_\allowbreak{}inverse.py}; the Gram certificate (including negative controls with a common root) is checked on 500 random polynomial families by \texttt{tools/verification/c316\_\allowbreak{}gram\_\allowbreak{}certificate\_\allowbreak{}random.py}; Cauchy--Binet and the shifted index are checked on 200 random pencils by \texttt{tools/verification/c317\_\allowbreak{}cauchy\_\allowbreak{}binet\_\allowbreak{}random.py}; and the linear system (R0) of Lemma 8.1 is compared with the gcd test on 135 random pencils, on common-root controls, on two constructed thin pencils whose minimal certificate degree is positive, and end-to-end on 51 zero-free-row instances (29 of them satisfying $(5a)$) by \texttt{tools/verification/c318\_\allowbreak{}r0\_\allowbreak{}linear\_\allowbreak{}system.py}. Lemma D2 and the data constraints (B.1)--(B.2) of the $\exists\mathbb R$ encoding are checked by \texttt{tools/verification/c319\_\allowbreak{}encoding\_\allowbreak{}and\_\allowbreak{}rowspace.py}: Lemma D2 on 40 random families, the encoding's witness direction on a constructed feasible instance (observable $(C,A)$, $L_0=e_1^{\mathsf T}$, $L=I_n$), and the necessity of (B.1)--(B.2) by perturbing $L_0$. These are finite exact checks, not proofs.


\begin{thebibliography}{99}
\bibitem{fer10} J. A. Fernando, H. Trinh, L. Jennings. Functional observability and the design of minimum order linear functional observers. \emph{IEEE Trans. Automat. Control} 55(5):1268--1273, 2010.
\bibitem{dar25} M. Darouach, J. A. Fernando. On functional observability and functional observer design. \emph{Automatica}, vol. 173, Art. no. 112115, 2025, doi: 10.1016/j.automatica.2025.112115.
\bibitem{fer26a} J. A. Fernando. A minimum-order functional observer beyond the Darouach and Luenberger constructions. arXiv:2609.33615, 2026 (preprint).
\bibitem{fer26b} J. A. Fernando. Complete characterization of minimum-order functional observers from Darouach to Luenberger. arXiv:2609.37992, 2026 (preprint).
\bibitem{ren92} J. Renegar. On the computational complexity and geometry of the first-order theory of the reals. Part II: The general decision problem. \emph{J. Symbolic Comput.} 13(3):301--327, 1992, doi: 10.1016/S0747-7171(10)80004-5.
\bibitem{bee88} Th. Beelen, P. Van Dooren. An improved algorithm for the computation of Kronecker's canonical form of a singular pencil. \emph{Linear Algebra Appl.} 105:9--65, 1988.
\bibitem{fab90} J.-P. Fabiano, D. C. Struppa. Degree bounds for polynomial matrix left inverses. \emph{J. Approx. Theory Appl.}, 1990.
\bibitem{emi99} I. Z. Emiris, B. Mourrain. Matrices in elimination theory. \emph{J. Symbolic Comput.} 28(1--2):3--44, 1999, doi: 10.1006/jsco.1998.0266.
\bibitem{nof23} V. Noferini, P. Van Dooren. On the stability of polynomial eigenvalue problems. \emph{Linear Algebra Appl.} 658, 2023.
\end{thebibliography}
\end{document}